\section{\label{I_2_C}Vers une refonte de S-Movies} 

Si le cadre posé par la modélisation FRBR s’avère limité pour des collections spécifiques telles que celles que nous avons analysé, le modèle induit des difficultés plus globales au niveau de l’interface et de la navigation utilisateur. Le FRBR est fondé sur des besoins utilisateurs dont l’objectif serait “d’identifier, sélectionner et acquérir” efficacement l’ensemble des éléments décrits et répartis entre différents nie=veaux hiérarchiues \footcite{zotero-item-646}. Ces besoins semblent respectés au sein de l’interface S-Movies, via la navigation permise entre les entités et les outils de recherche disponibles. Toutefois, une double complexité subsiste dans l’implémentation de la modélisation FRBR au sein d’une interface. D’une part, il reste complexe de mettre en place une interface intuitive permettant à l’utilisateur de connaître immédiatement le rôle descriptif de telle entité du FRBR et sa place au sein de la hiérarchie du modèle. D’autre part, la visualisation des liens entre les entités du FRBR reste limitée, puisqu’elle supposerait des développements supplémentaires innovants tels que la visualisation en graphe, par exemple.  

Parmi les complexités d’interface révélés par l'expérience utilisateur, certains éléments semblent anodins. Par exemple, l’affichage d’une hiérarchie inversée du FRBR dans le menu de l’application, pouvant nuire à l’intuitivité du rôle de chaque entité pour la description des collections. En effet, au sein du menu, les tables sont présentées de telle manière à ce que le plus bas niveau du FRBR (l'item) apparaisse en haut, suivi de la manifestation et du plus haut niveau (l'œuvre). Cet affichage visuel descendant peut induire en erreur l'utilisateur sur la navigation hiérarchique à employer au sein du modèle. De plus, au sein de S-Movies le plus bas niveau de la hiérarchie (l’Item) est désigné par la notion de “Film”. Or, dans un contexte cinématographique, le film peut communément désigner à la fois une œuvre et un item individuel. La notion de film n'évoque pas le caractère unique et individuel de l'exemplaire tel qu'il est défini par le manuel de la FIAF. Une autre confusion terminologique potentielle peut être créée par la présence de la notion d'Expression, alors même que ce niveau, issu du modèle FRBR bibliographique, est fusionné dans celui d'œuvre.  

La visualisation des liens entre les entités descriptives au sein de l’interface est également motrice d’améliorations au sein de l’espace de travail. L'interface actuelle permet, après configuration, de visualiser les liens d’une notice avec d’autres notices issues de niveaux différents du FRBR. Or, il serait pertinent pour l’utilisateur d’avoir la possibilité de visualiser les liens entre entités à un niveau plus large, et de manière graphique. Par exemple, il serait pertinent de pouvoir visualiser, au niveau d’une notice Œuvre, l'ensemble des manifestations dans lesquelles elle est matérialisée, ainsi que tous les exemplaires issus de ses manifestations respectives et conservés par l’institution. Cette visualisation des liens pourrait être étendue de manière plus large au niveau de la table d’une entité, ou de celui de la collection de l’institution.  

Par ailleurs, la mise en place de liens indissociables entre entités prévus par le modèle FRBR nécessitent, au sein de l’interface, des configurations spécifiques complexes. En effet, l’indissociabilité des liens entre les entités suppose que l’interface rende obligatoire pour l’utilisateur l’association de certaines notices entre elles. Par exemple, un Item doit obligatoirement être lié à une Manifestation et/ou à une Œuvre. Or l’obligation de ces liens n’est pas encore prévue par l’interface, laissant ainsi la possibilité à l’utilisateur de créer un Item sans Œuvre liée par exemple. De même, l’interface permet par exemple la suppression d'une notice Œuvre et de tous ses éléments liés (notices de Manifestations et d’Items), sans avertissement préalable. 

Ces difficultés d’interface sont en réalité issues de la complexité de traduire un cadre de modélisation défini, en termes d’expérience utilisateur. Cette traduction est cruciale afin d’assurer la qualité des données indexées au sein de l’interface. En effet, l’intuitivité de l'interface permet à l'utilisateur de diviser la description d’un objet entre les différents niveaux prévus par l’application, et d’associer ces niveaux entre eux afin qu’ils soient requêtables au sein de l’interface. L’absence d’intuitivité peut notamment empêcher l’héritage des données entre les différents niveaux, dans le cas d’une notice dépourvue de liens, par exemple. L’interface doit ainsi assurer la qualité des données en instaurant une circulation intuitive au sein du modèle FRBR.  